\chapter{Evidence provenance and supplementary inventory}
\label{app:provenance}

The numerical inputs are frozen exports of \filepath{thesis/evidence/}. The current export is taken from revision \texttt{\seqsplit{db028426c9a9f382f82b640b988ca9abdf313fd0}}; its evidence index records generation on 23 September 2026 at 12:04:28 UTC, with pipeline revision \texttt{\seqsplit{2b672b2cee025cb89408a4fcafc1d42edfd1ab99}}. The index warns that this pipeline revision need not be the revision used for every earlier measurement. Per-run records remain authoritative for historical boots. The raw physical-sweep source is recorded at 12:04:27 UTC on that date.

The previous fit is retained from revision \texttt{\seqsplit{5e1261632d8d89ee6878b3390cce092a261582be}} and the earlier evidence index generated on 22 September. Only its fit comparison is used as a historical comparison; the current physical tables use the refreshed export. Export revision, measurement identity, and document build date are distinct. The metadata record \filepath{evidence/origin.json} preserves the selected revisions without requiring a branch switch.

\Cref{tab:provenance} maps the figures, tables, and supporting findings to their inputs. In the submitted repository, \filepath{evidence/current/} corresponds to \filepath{thesis/evidence/} and \filepath{evidence/previous/} to \filepath{thesis/evidence/previous/}; other paths are relative to the repository root. Revision identifiers refer to the development history, which is not part of the submission; the SHA-256 values identify the submitted files. SHA-256 values are recomputed over the actual source-file bytes. They identify these documents and exports, not a new verification of remote SquashFS bytes. The frozen upstream provenance index retains raw source paths, dates, truncated upstream digests, and artefact freshness assessments.

\begingroup\footnotesize
\begin{longtable}{@{}p{3.1cm}p{5.4cm}p{4.1cm}@{}}
\caption{Source mapping for figures, tables, and supporting findings. Dates are evidence-export dates for frozen inputs and last-change dates at the pinned revision for repository documents; they are not all measurement dates. Digests identify source-file bytes.}\label{tab:provenance}\\
\toprule Figure / result & Source file and date & SHA-256 \\
\midrule\endfirsthead
\toprule Figure / result & Source file and date & SHA-256 \\
\midrule\endhead
\Cref{fig:timing,tab:physical} & \filepath{evidence/current/tier2-by-n.csv}\newline 2026-09-23 & \texttt{\seqsplit{fb8f80a68c82ed0d66a33c1e8ae6cb89c42f8dc420ee54f151949a132cc95066}} \\
\Cref{fig:timing,tab:fits,eq:combined-time} & \filepath{evidence/current/tier2-fit.csv}\newline 2026-09-23 & \texttt{\seqsplit{07d5f0ec00367d87d5a76b9735bee6bee0c6d8d78fc9ef0190099c4455eb6a05}} \\
Historical fit comparison & \filepath{evidence/previous/tier2-fit.csv}\newline 2026-09-22 & \texttt{\seqsplit{24067b86f18cb0c2cdec6bb11737f0e346f79a3a6824b53aaadc9564a86774b2}} \\
\Cref{tab:admission} & \filepath{evidence/current/tier1-totals.csv}\newline 2026-09-23 & \texttt{\seqsplit{b6698505f29ce5614a42576765cebefdd9dc5faaa9d3070691ceaa38e6a05308}} \\
\Cref{tab:causes} & \filepath{evidence/current/tier1-classes.csv}\newline 2026-09-23 & \texttt{\seqsplit{4bca7c7ce81d15da826e2bad5fa49e38b794faee2799d13d6e9e18b951f7c70a}} \\
Control and set extension & \filepath{evidence/current/tier1-control.csv}\newline 2026-09-23 & \texttt{\seqsplit{34fd65f843aed9e8a7ec9a824d4222fbbc0ac84e89eb49513c6558c2586193da}} \\
Set-extension cross-check & \filepath{evidence/current/tier1-crosscheck.csv}\newline 2026-09-23 & \texttt{\seqsplit{4311d2f149cb11a871e86e285be4dfe92ccc42550cbc6992b998f89a45a906fd}} \\
Non-monotone exceptions & \filepath{evidence/current/tier1-nonmonotone.csv}\newline 2026-09-23 & \texttt{\seqsplit{69f233ae068b672506c52786871181262c7ebd19ef7813fa9c3560d6d698eb83}} \\
Permitted overlap & \filepath{evidence/current/tier1-measured.csv}\newline 2026-09-23 & \texttt{\seqsplit{97de92fcf54cb0088b4cc12f8ecd7cef7468bca4bbd848e1163a2b1c975b2bb3}} \\
V1--V8 outcomes & \filepath{evidence/current/tier2-checks.csv}\newline 2026-09-23 & \texttt{\seqsplit{fde1c2697cd986acc202c77306a5ebebcb72d4d397803728b6d69cdaff8159d1}} \\
\Cref{fig:storage,tab:storage,eq:storage} & \filepath{evidence/current/storage-cohorts.csv}\newline 2026-09-23 & \texttt{\seqsplit{b5b065efe366840b063c25ad880df31e6a8360d7f883e26c055b31dcf673d81e}} \\
Storage exclusions & \filepath{evidence/current/storage-sensitivity.csv}\newline 2026-09-23 & \texttt{\seqsplit{a8a62be93a979cc0a231f5ef8fb0f2f589883810e53fb44f6710a4397272e64b}} \\
\Cref{tab:calibration} & \filepath{evidence/current/storage-model-check.csv}\newline 2026-09-23 & \texttt{\seqsplit{d1a770f0038d7cfb65d286c1afa8c7a0071686428b3da944998b7095cc39cdae}} \\
\Cref{tab:module-inventory} & \filepath{evidence/current/storage-per-module.csv}\newline 2026-09-23 & \texttt{\seqsplit{eca500b382a74f1ea7c54d440127f4ae6506508e6ecd5a1d822a9051e426d595}} \\
\Cref{tab:boots}; boot durations & \filepath{evidence/current/tier3.csv}\newline 2026-09-23 & \texttt{\seqsplit{e5edb5f45b5f878e31111003b71b8a6a5b01047be88268650c31922d4fbc3ead}} \\
Catalogue classification & \filepath{evidence/current/catalogue.csv}\newline 2026-09-23 & \texttt{\seqsplit{ba3951d6219fcade0283d9b41e21d5e8fa91c53be565d6d0dc336ca8fdeee706}} \\
Historical fat-base model; timing exclusions & \filepath{evidence/current/claims.md}\newline 2026-09-23 & \texttt{\seqsplit{b97ddf10d318d7257000b907ad40fb515bf4d227698d509acace1c04f13445f3}} \\
Evidence corrections & \filepath{evidence/current/inconsistencies.md}\newline 2026-09-23 & \texttt{\seqsplit{477c82e758c77021c2189037abef530f2c7a17ae2d716bed1747a6816ff95771}} \\
Upstream identities and dates & \filepath{evidence/current/provenance.md}\newline 2026-09-23 & \texttt{\seqsplit{6f797c62a1eb685a00d8974ac572700a370c8bd25c83051076485917ed2798e4}} \\
Previous upstream identities & \filepath{evidence/previous/provenance.md}\newline 2026-09-22 & \texttt{\seqsplit{0d3ccd27ebcbcdd242f8bbc2edd6cfb77a6fba30c5dd2b0d2aab3da7f0223b8c}} \\
Base bytes and recorded artefact identities & \filepath{evidence/current/provenance-artefacts.csv}\newline 2026-09-23 & \texttt{\seqsplit{42fb744be2ccc74590612672c2cfd1646b801f0da9df1ddce3e3e33431aaa82b}} \\
\Cref{fig:architecture,tab:taxonomy,tab:digests} & \filepath{ARCHITECTURE.md}\newline 2026-09-24 & \texttt{\seqsplit{9ab8d118c42c5b038708d8ed39fb691e7ea97497c87b29dc0d8b573adcaa080f}} \\
\Cref{tab:manifest} & \filepath{docs/DATA_MODEL.md}\newline 2026-09-23 & \texttt{\seqsplit{8d1fcefa209f038e49a11eb29fea63898001e0122650b25e16b870bb6ebca5a1}} \\
Registry discovery; schema history & \filepath{ARCHITECTURE.md}\newline 2026-09-24 & \texttt{\seqsplit{9ab8d118c42c5b038708d8ed39fb691e7ea97497c87b29dc0d8b573adcaa080f}} \\
Physical sampling and guest observation & \filepath{docs/SAMPLING_AND_BOOT.md}\newline 2026-09-23 & \texttt{\seqsplit{287ce7d675aa3d9180f91ebf8511a0530fbdd0e0337dc0b7bb2462b13ed81520}} \\
Merger counterexamples & \filepath{tests/round4_attacks.py}\newline 2026-09-23 & \texttt{\seqsplit{a5930cf1376c680dbab61a39d307abe2f3d7138ed6e2c25d417ddacf348380ff}} \\
Equal-size substitution & \filepath{tests/round2_attacks.py}\newline 2026-09-19 & \texttt{\seqsplit{a0773989399c7ded6597d29ef26efc3cf2c72bec41e0828487912ee66419ff41}} \\
Evidence-pipeline counterexamples & \filepath{tests/round3_evidence.py}\newline 2026-09-21 & \texttt{\seqsplit{7ded29d4a41010e645aaefc24e736cbd51ce68c72ce9756e6da2e5d02fe21a04}} \\
Registry implementation & \filepath{scripts/reconcile.py}\newline 2026-09-23 & \texttt{\seqsplit{24b4f3bfb591c63ddad3ea21e899b27e3412be4abaae669defef02abd7df871c}} \\
Registry formats and implementation boundaries & \filepath{docs/REGISTRY_FORMATS.md}\newline 2026-09-23 & \texttt{\seqsplit{a6c48c5a20559ed3766c2d90403510baca32d63c863b181f44d27d5f55fe3e10}} \\
Independent monolith construction & \filepath{scripts/02_build_delta.sh}\newline 2026-09-20 & \texttt{\seqsplit{d28ff1db2bb0aef9ac96eac16c2f1cbe5900a519e4c28150c7cd292e17fb16c0}} \\
Separate flattened-root comparator & \filepath{scripts/15_measure_monolith.sh}\newline 2026-09-19 & \texttt{\seqsplit{afc8c7972bafbe6a67491f9324e84bfd10a8cdaca346fff8afc6fab3e1991ccf}} \\
Independent rebuild mechanisms and userspace GPU test & \filepath{docs/evidence/remote-verification-2026-09-20/REPORT.md}\newline 2026-09-20 & \texttt{\seqsplit{0417c49b8705eb53bba6dadde2fb7ffa9fff0411a9cf58b87c738b0871c7052e}} \\
Declaration and disclosure requirements & \filepath{thesis/thesis guidelines/AI_usage_in_thesis_work.pdf}\newline 2026-09-20 & \texttt{\seqsplit{fbf2bdde6d980f485e8bf81a8b1706006eb75dcaa7c31f929eafcb37a57f10a1}} \\
\bottomrule\end{longtable}\endgroup


\section{Derivations and evidence qualifications}

\filepath{prepare_data.py} generates \filepath{data.tex}, the numerical tables, and vector plot coordinates. It reads frozen exports and performs no experiments. The combined timing fit sums the current compose and verify coefficients; it assigns no combined coefficient of determination. The slope and intercept comparisons use the previous composition fit only and make no causal attribution to hardware or algorithm.

The older quoted Tier-2-to-boot interval of \numrange{62}{686} used the superseded physical fit. The chapter recomputes descriptive comparisons by layer count using the refreshed combined fit and the retained tier-3 run durations. Neither interval is a matched experiment across identical inputs. Admission timings marked untraceable in the evidence are not used for a cross-tier ratio.

The storage numerator is the additive model for every module, including those with a measured comparator. The measured monoliths calibrate that model; they are not substituted into the headline sum. The historical fat-base values are traced through the current claims index and remain labelled as a different catalogue. The derived percentage is computed from those rounded historical totals.

The per-composition rows of the physical sweep are included as \filepath{thesis/evidence/raw/compose-sweep.csv}, so its fits can be recomputed; the file's SHA-256 matches the source digest recorded in the exported provenance index. Blank verdict cells in the boot CSV mean no result. They are counted separately from explicit ABORTED and BROKEN verdicts. Rows classified PASS or FAIL supply the completed-duration range. The selected rows in the boot table retain their original order and identities in the frozen CSV; no result is inferred for unfinished runs.

\section{Evidence not collected}

Four measurements that would strengthen the evaluation were not made.

\begin{itemize}
\item No quantitative summary of the independent rebuild was exported into the thesis evidence. The body therefore describes the cross-machine differences and their mechanisms without stating a reproduction rate.
\item No matched monolith was built for the large workloads. The monolithic comparison is calibrated on small modules only, and the headline storage saving remains modelled.
\item Timings were not repeated under controlled conditions. The timing figure shows means and fits without dispersion or uncertainty bands, and it supports no causal cross-machine performance claim.
\item The artefacts of the final composition sweep were not booted. The reported boot results come from earlier artefacts and are labelled as such.
\end{itemize}


\section{Complete compressed module inventory}

\Cref{tab:module-inventory} provides the individual values behind the storage discussion. The large CUDA delta explains much of the difference between small-cohort and whole-catalogue sharing ratios. The marginal saving column assumes a stored base; it is not the saving of an isolated one-module repository. The control and synthetic entries are retained because the headline inventory includes them.

\begin{longtable}{lrrr}
\caption{Compressed module inventory for the \dref{/catalogue/total}-module catalogue. Marginal saving assumes the base is already stored.}\label{tab:module-inventory}\\
\toprule Module & Delta (\si{\mega\byte}) & Model (\si{\mega\byte}) & Saving (\si{\percent}) \\
\midrule\endfirsthead
\toprule Module & Delta (\si{\mega\byte}) & Model (\si{\mega\byte}) & Saving (\si{\percent}) \\
\midrule\endhead
cuda-runtime & \dref{/module/cuda-runtime/delta-mb} & \dref{/module/cuda-runtime/monolithic-mb} & \dref{/module/cuda-runtime/saving-pct} \\
nvidia-driver-535 & \dref{/module/nvidia-driver-535/delta-mb} & \dref{/module/nvidia-driver-535/monolithic-mb} & \dref{/module/nvidia-driver-535/saving-pct} \\
rust & \dref{/module/rust/delta-mb} & \dref{/module/rust/monolithic-mb} & \dref{/module/rust/saving-pct} \\
java & \dref{/module/java/delta-mb} & \dref{/module/java/monolithic-mb} & \dref{/module/java/saving-pct} \\
llvm & \dref{/module/llvm/delta-mb} & \dref{/module/llvm/monolithic-mb} & \dref{/module/llvm/saving-pct} \\
gcc & \dref{/module/gcc/delta-mb} & \dref{/module/gcc/monolithic-mb} & \dref{/module/gcc/saving-pct} \\
docker & \dref{/module/docker/delta-mb} & \dref{/module/docker/monolithic-mb} & \dref{/module/docker/saving-pct} \\
postgres & \dref{/module/postgres/delta-mb} & \dref{/module/postgres/monolithic-mb} & \dref{/module/postgres/saving-pct} \\
mysql & \dref{/module/mysql/delta-mb} & \dref{/module/mysql/monolithic-mb} & \dref{/module/mysql/saving-pct} \\
emacs & \dref{/module/emacs/delta-mb} & \dref{/module/emacs/monolithic-mb} & \dref{/module/emacs/saving-pct} \\
apache & \dref{/module/apache/delta-mb} & \dref{/module/apache/monolithic-mb} & \dref{/module/apache/saving-pct} \\
pytools & \dref{/module/pytools/delta-mb} & \dref{/module/pytools/monolithic-mb} & \dref{/module/pytools/saving-pct} \\
webserver & \dref{/module/webserver/delta-mb} & \dref{/module/webserver/monolithic-mb} & \dref{/module/webserver/saving-pct} \\
vim & \dref{/module/vim/delta-mb} & \dref{/module/vim/monolithic-mb} & \dref{/module/vim/saving-pct} \\
git & \dref{/module/git/delta-mb} & \dref{/module/git/monolithic-mb} & \dref{/module/git/saving-pct} \\
dnsutils & \dref{/module/dnsutils/delta-mb} & \dref{/module/dnsutils/monolithic-mb} & \dref{/module/dnsutils/saving-pct} \\
pipdemo & \dref{/module/pipdemo/delta-mb} & \dref{/module/pipdemo/monolithic-mb} & \dref{/module/pipdemo/saving-pct} \\
pgclient & \dref{/module/pgclient/delta-mb} & \dref{/module/pgclient/monolithic-mb} & \dref{/module/pgclient/saving-pct} \\
jq & \dref{/module/jq/delta-mb} & \dref{/module/jq/monolithic-mb} & \dref{/module/jq/saving-pct} \\
memcached & \dref{/module/memcached/delta-mb} & \dref{/module/memcached/monolithic-mb} & \dref{/module/memcached/saving-pct} \\
pyyaml & \dref{/module/pyyaml/delta-mb} & \dref{/module/pyyaml/monolithic-mb} & \dref{/module/pyyaml/saving-pct} \\
gawk & \dref{/module/gawk/delta-mb} & \dref{/module/gawk/monolithic-mb} & \dref{/module/gawk/saving-pct} \\
mta-msmtp & \dref{/module/mta-msmtp/delta-mb} & \dref{/module/mta-msmtp/monolithic-mb} & \dref{/module/mta-msmtp/saving-pct} \\
sqlite & \dref{/module/sqlite/delta-mb} & \dref{/module/sqlite/monolithic-mb} & \dref{/module/sqlite/saving-pct} \\
redis & \dref{/module/redis/delta-mb} & \dref{/module/redis/monolithic-mb} & \dref{/module/redis/saving-pct} \\
control-oldsnap & \dref{/module/control-oldsnap/delta-mb} & \dref{/module/control-oldsnap/monolithic-mb} & \dref{/module/control-oldsnap/saving-pct} \\
curl & \dref{/module/curl/delta-mb} & \dref{/module/curl/monolithic-mb} & \dref{/module/curl/saving-pct} \\
tcpdump & \dref{/module/tcpdump/delta-mb} & \dref{/module/tcpdump/monolithic-mb} & \dref{/module/tcpdump/saving-pct} \\
zstd & \dref{/module/zstd/delta-mb} & \dref{/module/zstd/monolithic-mb} & \dref{/module/zstd/saving-pct} \\
tmux & \dref{/module/tmux/delta-mb} & \dref{/module/tmux/monolithic-mb} & \dref{/module/tmux/saving-pct} \\
rsync & \dref{/module/rsync/delta-mb} & \dref{/module/rsync/monolithic-mb} & \dref{/module/rsync/saving-pct} \\
socat & \dref{/module/socat/delta-mb} & \dref{/module/socat/monolithic-mb} & \dref{/module/socat/saving-pct} \\
wget & \dref{/module/wget/delta-mb} & \dref{/module/wget/monolithic-mb} & \dref{/module/wget/saving-pct} \\
mta-nullmailer & \dref{/module/mta-nullmailer/delta-mb} & \dref{/module/mta-nullmailer/monolithic-mb} & \dref{/module/mta-nullmailer/saving-pct} \\
htop & \dref{/module/htop/delta-mb} & \dref{/module/htop/monolithic-mb} & \dref{/module/htop/saving-pct} \\
nc-openbsd & \dref{/module/nc-openbsd/delta-mb} & \dref{/module/nc-openbsd/monolithic-mb} & \dref{/module/nc-openbsd/saving-pct} \\
original-awk & \dref{/module/original-awk/delta-mb} & \dref{/module/original-awk/monolithic-mb} & \dref{/module/original-awk/saving-pct} \\
nc-traditional & \dref{/module/nc-traditional/delta-mb} & \dref{/module/nc-traditional/monolithic-mb} & \dref{/module/nc-traditional/saving-pct} \\
fake-nvidia-driver & \dref{/module/fake-nvidia-driver/delta-mb} & \dref{/module/fake-nvidia-driver/monolithic-mb} & \dref{/module/fake-nvidia-driver/saving-pct} \\
fake-cuda & \dref{/module/fake-cuda/delta-mb} & \dref{/module/fake-cuda/monolithic-mb} & \dref{/module/fake-cuda/saving-pct} \\
\bottomrule\end{longtable}

